% ============================================================================
% MUSTERLÖSUNG: Elektromagnetische Induktion (LP-09b)
% Physik Klasse 9
% Stand: 03.02.2026
% ============================================================================

\documentclass[11pt,a4paper]{article}

\usepackage[utf8]{inputenc}
\usepackage[T1]{fontenc}
\usepackage[ngerman]{babel}
\usepackage{geometry}
\usepackage{helvet}
\usepackage{enumitem}
\usepackage{fancyhdr}
\usepackage{xcolor}
\usepackage{tcolorbox}
\usepackage{amssymb}
\usepackage{siunitx}
\usepackage{array}
\usepackage{textcomp}
\usepackage{qrcode}
\usepackage[hidelinks]{hyperref}

\sisetup{locale=DE, per-mode=symbol}
\renewcommand{\familydefault}{\sfdefault}

\geometry{left=1.4cm, right=1.4cm, top=1.3cm, bottom=1.0cm}

% Fachfarbe Physik
\definecolor{physik}{HTML}{2c5aa0}
\colorlet{fachfarbe}{physik}
\definecolor{protokollrand}{HTML}{2a7a4b}
\definecolor{merkerand}{HTML}{b35c00}
\definecolor{loesung}{HTML}{c62828}

% Kopfzeile
\pagestyle{fancy}
\fancyhf{}
\fancyhead[L]{\small\textbf{Physik Klasse 9} | Elektromagnetische Induktion}
\fancyhead[R]{\small\textcolor{loesung}{\textbf{MUSTERLÖSUNG}}}
\renewcommand{\headrulewidth}{0.4pt}
\setlength{\headheight}{14pt}

% Boxen
\tcbuselibrary{skins}

\newtcolorbox{merkkasten}[1][]{
    colback=white, colframe=fachfarbe,
    coltitle=black, colbacktitle=white,
    fonttitle=\bfseries\small,
    boxrule=1pt, arc=0pt,
    left=6pt, right=6pt, top=4pt, bottom=4pt,
    #1
}

\newtcolorbox{protokollbox}[1][]{
    colback=white, colframe=protokollrand,
    coltitle=black, colbacktitle=white,
    fonttitle=\bfseries\small,
    boxrule=1pt, arc=0pt,
    left=6pt, right=6pt, top=4pt, bottom=4pt,
    #1
}

\newtcolorbox{merkebox}[1][]{
    colback=white, colframe=merkerand,
    coltitle=black, colbacktitle=white,
    fonttitle=\bfseries\small,
    boxrule=1pt, arc=0pt,
    left=6pt, right=6pt, top=4pt, bottom=4pt,
    #1
}

\newcommand{\aufgabe}[1]{\noindent\textbf{\textcolor{fachfarbe}{#1}}}
\newcommand{\lsg}[1]{\textcolor{loesung}{\textbf{#1}}}

\setlength{\parindent}{0pt}
\setlength{\parskip}{0pt}

\begin{document}

% ============================================================================
% TITEL
% ============================================================================
\begin{minipage}[t]{0.80\textwidth}
    \vspace{0pt}
    {\Large\bfseries Elektromagnetische Induktion}\\[0.2em]
    {\small\textcolor{loesung}{\textbf{MUSTERLÖSUNG}} \hfill Datum: 03.02.2026}
\end{minipage}
\vspace{0.1em}
\hrule
\vspace{0.4em}

% ============================================================================
% LEITFRAGE
% ============================================================================
\begin{merkkasten}
\textbf{Leitfrage:} Der Motor macht aus Strom Bewegung. \textbf{Geht das auch umgekehrt?}
\end{merkkasten}

\vspace{0.3em}

% ============================================================================
% KASTEN 1
% ============================================================================
\begin{protokollbox}[title=Kasten 1: Simulationsexperiment -- Protokoll]
\small
\renewcommand{\arraystretch}{1.25}
\begin{tabular}{|>{\raggedright\arraybackslash}p{4.2cm}|p{2.2cm}|p{7.5cm}|}
\hline
\textbf{Experiment} & \textbf{Spannung $U$} & \textbf{Beobachtung} \\
\hline
E1: Magnet langsam durch Spule & \lsg{ca. 0,5 V} & \lsg{Zeiger schlägt leicht aus} \\[0.15cm]
\hline
E2: Magnet schnell durch Spule & \lsg{ca. 2--3 V} & \lsg{Zeiger schlägt stark aus} \\[0.15cm]
\hline
E3: Magnet ruht in der Spule & \lsg{0 V} & \lsg{Zeiger bleibt auf Null} \\[0.15cm]
\hline
E4: Bewegungsrichtung umkehren & \lsg{negativ} & \lsg{Zeiger schlägt in andere Richtung aus} \\[0.15cm]
\hline
\end{tabular}

\vspace{0.3em}
\textbf{Schlussfolgerung:} Eine Spannung entsteht nur, wenn \lsg{sich der Magnet bewegt / eine Änderung stattfindet}.
\end{protokollbox}

\vspace{0.3em}

% ============================================================================
% KASTEN 2
% ============================================================================
\begin{protokollbox}[title=Kasten 2: Einflussfaktoren -- Was erhöht die Spannung?]
\small
\renewcommand{\arraystretch}{1.2}
\begin{tabular}{|p{3.2cm}|c|c|c|p{4cm}|}
\hline
\textbf{Variable} & \textbf{Einst. 1} & \textbf{Einst. 2} & \textbf{$U$ größer?} & \textbf{Erkenntnis} \\
\hline
Geschwindigkeit & langsam & schnell & \lsg{$\boxtimes$ Ja} ~ $\square$ Nein & \lsg{schneller = größere Änd.} \\[0.1cm]
\hline
Windungszahl n & 10 & 50 & \lsg{$\boxtimes$ Ja} ~ $\square$ Nein & \lsg{5$\times$ n $\approx$ 5$\times$ U} \\[0.1cm]
\hline
Magnetstärke & normal & stark & \lsg{$\boxtimes$ Ja} ~ $\square$ Nein & \lsg{mehr Feldlinien} \\[0.1cm]
\hline
Spulenfläche A & klein & groß & \lsg{$\boxtimes$ Ja} ~ $\square$ Nein & \lsg{mehr FL durchsetzen} \\[0.1cm]
\hline
\end{tabular}

\vspace{0.2em}
\textbf{Fazit:} Die Induktionsspannung wird größer durch:
1. \lsg{schnellere Bew.} ~ 2. \lsg{mehr Windungen} ~ 3. \lsg{stärkeren Magn.} ~ 4. \lsg{größere Fläche}
\end{protokollbox}

\vspace{0.3em}

% ============================================================================
% KASTEN 3
% ============================================================================
\begin{merkebox}[title=Kasten 3: Merksatz -- Elektromagnetische Induktion]
\small
\textbf{Definition:} Bewegt man einen \lsg{Magneten} relativ zu einer \lsg{Spule}, so entsteht eine \lsg{Induktionsspannung}. Diesen Vorgang nennt man \textbf{Induktion}.

\vspace{0.3em}
\textbf{Voraussetzung:} Es muss eine \lsg{Bewegung / Änderung} geben -- ohne Bewegung keine Spannung!

\vspace{0.3em}
\textbf{Kernidee:} Eine Spannung entsteht, wenn sich die Anzahl der \lsg{Feldlinien}, die die Spule durchsetzen, \textbf{ändert}.

\vspace{0.3em}
\textbf{Induktionsgesetz (qualitativ):} Die Induktionsspannung $U_{\text{ind}}$ ist umso größer, je ...
\begin{itemize}[leftmargin=1.2cm, itemsep=0pt, topsep=2pt]
    \item[$\triangleright$] \lsg{schneller} die Bewegung ist (= schnellere Änderung).
    \item[$\triangleright$] \lsg{stärker} der Magnet ist (= mehr Feldlinien).
    \item[$\triangleright$] \lsg{mehr Windungen} die Spule hat.
    \item[$\triangleright$] \lsg{größer} die Spulenfläche ist (= mehr Feldlinien durchsetzen).
\end{itemize}
\end{merkebox}

% ============================================================================
% SEITE 2
% ============================================================================
\newpage

% ============================================================================
% KASTEN 4
% ============================================================================
\begin{merkkasten}
\textbf{Kasten 4: Vergleich Motor und Generator}

\vspace{0.2em}
\small
Motor und Generator haben den \textbf{gleichen Aufbau}, aber \textbf{umgekehrte Funktionen}.

\vspace{0.3em}
\renewcommand{\arraystretch}{1.3}
\begin{tabular}{|p{3.5cm}|p{4.8cm}|p{5cm}|}
\hline
& \textbf{Motor} & \textbf{Generator} \\
\hline
Energieumwandlung & Elektr. Energie $\rightarrow$ \lsg{Bewegung} & \lsg{Bewegung} $\rightarrow$ Elektr. Energie \\
\hline
Strom & fließt \lsg{hinein} & kommt \lsg{heraus} \\
\hline
Funktionsweise & Strom rein $\rightarrow$ Bewegung & \lsg{Bewegung $\rightarrow$ Spannung raus} \\
\hline
Beispiel & E-Bike, Ventilator & \lsg{Kraftwerk, Dynamo} \\
\hline
\end{tabular}

\vspace{0.3em}
\textbf{Merksatz:} Motor und Generator sind „\lsg{Zwillinge / Umkehrungen}" -- gleicher Aufbau, umgekehrte Funktion!
\end{merkkasten}

\vspace{0.4em}

% ============================================================================
% ÜBUNGEN
% ============================================================================
\aufgabe{Übung 1: Anwendungen der Induktion}

\vspace{0.2em}
{\small
\renewcommand{\arraystretch}{1.2}
\begin{tabular}{|p{3cm}|p{11cm}|}
\hline
\textbf{Anwendung} & \textbf{Wie funktioniert die Induktion hier?} \\
\hline
a) Fahrraddynamo & \lsg{Das Rad dreht einen Magneten neben einer Spule. Durch die Drehung ändert sich das Magnetfeld in der Spule ständig $\rightarrow$ Induktionsspannung $\rightarrow$ Strom für die Lampe.} \\[0.1cm]
\hline
b) Wasserkraftwerk & \lsg{Fallendes Wasser treibt eine Turbine an, die einen Generator dreht. Im Generator rotiert eine Spule im Magnetfeld $\rightarrow$ Induktion $\rightarrow$ Wechselstrom.} \\[0.1cm]
\hline
c) Windkraftanlage & \lsg{Der Wind dreht die Rotorblätter, die einen Generator antreiben. Im Generator wird durch Induktion elektrische Energie erzeugt.} \\[0.1cm]
\hline
\end{tabular}
}

\vspace{0.4em}

\aufgabe{Übung 2: Richtig oder falsch?}

\vspace{0.2em}
{\small
\begin{tabular}{|p{12cm}|c|c|}
\hline
\textbf{Aussage} & \textbf{R} & \textbf{F} \\
\hline
a) Ein ruhender Magnet in einer Spule erzeugt eine Spannung. & $\square$ & \lsg{$\boxtimes$} \\
\hline
b) Je schneller die Bewegung, desto größer die Induktionsspannung. & \lsg{$\boxtimes$} & $\square$ \\
\hline
c) Der Generator wandelt Bewegungsenergie in elektrische Energie um. & \lsg{$\boxtimes$} & $\square$ \\
\hline
d) Im Kraftwerk werden Motoren zur Stromerzeugung verwendet. & $\square$ & \lsg{$\boxtimes$} \\
\hline
e) Richtungswechsel der Bewegung ändert das Vorzeichen der Spannung. & \lsg{$\boxtimes$} & $\square$ \\
\hline
\end{tabular}
}

\vspace{0.4em}

\aufgabe{Übung 3: Transferaufgabe}

\vspace{0.2em}
{\small Eine Schülerin schüttelt eine Faraday-Taschenlampe (ohne Batterien) und die Lampe leuchtet.}

\vspace{0.2em}
a) Erkläre, wie die Lampe funktioniert. Nutze den Begriff \textit{Induktion}.

\vspace{0.1em}
\lsg{In der Lampe befindet sich ein Magnet, der durch eine Spule gleitet. Beim Schütteln bewegt sich der Magnet hin und her durch die Spule. Durch Induktion entsteht eine Spannung, die einen Kondensator auflädt. Der Kondensator versorgt die LED mit Strom.}

\vspace{0.3em}
b) Warum leuchtet die Lampe heller, wenn man schneller schüttelt?

\vspace{0.1em}
\lsg{Schnelleres Schütteln bedeutet schnellere Bewegung des Magneten. Je schneller die Bewegung, desto größer die Induktionsspannung. Größere Spannung $\rightarrow$ mehr Strom $\rightarrow$ helleres Leuchten.}

\end{document}
